\documentclass[11pt,a4paper]{article}
% =====================================================================
%  Gemeinsame Vorlage der Arbeitsblätter
%  "Verschlüsselung" – Informatik 11 (NTG)
%  (Antwortplatz-Makros wie informatik/10/java/arbeitsblaetter/ab-vorlage.tex)
%  Einbinden in jedem Blatt direkt nach \documentclass: % =====================================================================
%  Gemeinsame Vorlage der Arbeitsblätter
%  "Verschlüsselung" – Informatik 11 (NTG)
%  (Antwortplatz-Makros wie informatik/10/java/arbeitsblaetter/ab-vorlage.tex)
%  Einbinden in jedem Blatt direkt nach \documentclass: % =====================================================================
%  Gemeinsame Vorlage der Arbeitsblätter
%  "Verschlüsselung" – Informatik 11 (NTG)
%  (Antwortplatz-Makros wie informatik/10/java/arbeitsblaetter/ab-vorlage.tex)
%  Einbinden in jedem Blatt direkt nach \documentclass: \input{ab-vorlage}
%  Übersetzen mit XeLaTeX, LuaLaTeX, pdfLaTeX oder Tectonic.
%
%  Blatt:   tectonic ab-5-1-grundlagen.tex   (oder xelatex)
%  Lösung:  tectonic ab-5-1-grundlagen-lsg.tex
%           (Datei enthält nur: \def\mitloesung{}\input{ab-5-1-grundlagen};
%            füllt alle Lücken rot aus; siehe \lsg, \lsgrest, \lsglinien,
%            \lsgfeld, \lsgkariert, \lsgkarobox, \lsgtext, \lsgoder)
%  Lösungseinträge belegen nie mehr Platz als die Lücke im Blatt.
% =====================================================================
\usepackage[a4paper,left=18mm,right=18mm,top=26mm,bottom=17mm,headheight=15mm,headsep=4mm,footskip=8mm]{geometry}
\usepackage{iftex}
\ifPDFTeX
  \usepackage[T1]{fontenc}
  \usepackage[utf8]{inputenc}
  \usepackage{helvet}
  \usepackage[scaled=0.86]{beramono}
  \renewcommand{\familydefault}{\sfdefault}
\else
  \usepackage{fontspec}
  \setmainfont{texgyreheros}[Extension=.otf,UprightFont=*-regular,BoldFont=*-bold,ItalicFont=*-italic,BoldItalicFont=*-bolditalic]
  \setmonofont{DejaVuSansMono}[Extension=.ttf,UprightFont=*,BoldFont=*-Bold,ItalicFont=*-Oblique,BoldItalicFont=*-BoldOblique,Scale=0.86]
\fi
\usepackage[ngerman,shorthands=off]{babel}
\usepackage{xcolor,graphicx,tabularx,array,enumitem,fancyhdr,tikz,listings,multicol,calc,etoolbox,ragged2e,colortbl,xspace,pifont}
\usepackage[most]{tcolorbox}
\usetikzlibrary{arrows.meta,positioning,calc,decorations.pathreplacing,shapes.geometric,shapes.symbols,fit}
\usepackage[hidelinks]{hyperref}
\setlength{\parindent}{0pt}
\setlength{\parskip}{3pt}
\setlist{nosep,leftmargin=*}

% ---------- Lösungsschalter ----------
\newif\ifloesung
\ifdefined\mitloesung\loesungtrue\fi
\newcommand{\lsgfarbe}{\color{red!75!black}}

% ---------- Kopf- und Fußzeile (einheitliche Form) ----------
\newcommand{\lehrkraft}{Hau}
\newcommand{\themenbereich}{5. Verschlüsselung}
\newcommand{\abnummer}{}
\pagestyle{fancy}\fancyhf{}
\renewcommand{\headrulewidth}{0.4pt}
\fancyhead[L]{\small Klasse 11\\Datum:}
\fancyhead[C]{\small Themenbereich:\\\themenbereich}
\fancyhead[R]{\small \lehrkraft\ifloesung\\{\lsgfarbe\bfseries Lösung}\fi}
\fancyfoot[C]{\scriptsize edu-mrh.de\,·\,Informatik 11\,·\,Blatt \abnummer\ifloesung\ (Lösung)\fi\,·\,Seite \thepage}
\newcommand{\abtitel}[2]{\renewcommand{\abnummer}{#1}\begin{center}\Large\bfseries #1\quad\underline{#2}\end{center}\vspace{-1mm}}

% ---------- Boxen ----------
\newenvironment{aufgabe}[2][]{\begin{tcolorbox}[enhanced,breakable,colback=white,colframe=black,boxrule=0.7pt,arc=3mm,left=2mm,right=2mm,top=1.5mm,bottom=1.5mm,before skip=2.5mm,after skip=2.5mm]\textbf{Aufgabe #2}\ifstrempty{#1}{}{ \textit{#1}}\par\vspace{0.6mm}}{\end{tcolorbox}}
\newtcolorbox{merke}[1][Merke]{enhanced,breakable,colback=green!5,colframe=green!40!black,boxrule=0.8pt,arc=2mm,left=2mm,right=2mm,top=1.5mm,bottom=1.5mm,before skip=2.5mm,after skip=2.5mm,title={#1},fonttitle=\bfseries,coltitle=white}
\newtcolorbox{hinweis}{enhanced,breakable,colback=yellow!12,colframe=orange!70!black,boxrule=0.6pt,arc=2mm,left=2mm,right=2mm,top=1mm,bottom=1mm,before skip=2mm,after skip=2mm}
\newcommand{\web}[1]{{\small\color{gray}\textit{Website: #1}}}
% Kasten für optionale Aufgaben
\newtcolorbox{optional}[1][Optional]{enhanced,breakable,colback=blue!3,colframe=blue!45!black,boxrule=0.6pt,arc=2mm,left=2mm,right=2mm,top=1.5mm,bottom=1.5mm,before skip=2.5mm,after skip=2.5mm,title={#1},fonttitle=\bfseries\small,coltitle=white}
% Zahl mit Basis: \zs{1011}{2}
\newcommand{\zs}[2]{#1\textsubscript{#2}}

% ---------- Lücken, Linien, Karos ----------
% Einheitliche Maße für alle Blätter der Sequenz:
%  \linien{n}       n Schreiblinien, Abstand 8 mm, erste Linie 8 mm unter dem Text
%  \kariert{n}      Karoraster 5 mm, n Kästchen hoch, nur ganze Kästchen, mittig
%  \karobox{b}{h}   Karoraster b x h, auf ganze Kästchen gerundet
%  \luecke[l]       Lücke im Fließtext (Standard 4,5 cm, Fachbegriffe mind. 3,5 cm)
%  (jeweils mit Lösungsvariante \lsg…, siehe Kopf der Datei)
%  \restzeile       Lücke bis zum Zeilenende
%  \lueckentext     am Anfang eines Absatzes/Kastens mit Lücken: größerer
%                   Zeilenabstand und Flattersatz, damit man in die Lücken schreiben kann
%  \linienfeld{n}   umrahmtes Feld mit n Schreiblinien (z. B. für Java-Code)
%  \schreibzeile    Strut für Tabellenzellen, in die geschrieben wird (8 mm hoch)
\newlength{\zeilenabstand}\setlength{\zeilenabstand}{8mm}
\newlength{\karo}\setlength{\karo}{5mm}
\tikzset{karo/.style={draw=gray!50,line width=0.3pt},schreiblinie/.style={draw=black!65,line width=0.35pt}}
\newlength{\lsgbreite}
% Lücke: \luecke[Breite] / mit Lösung \lsg[Breite]{Text} (zu langer Text wird verkleinert)
\newcommand{\lsg}[2][4.5cm]{\leavevmode\ifloesung\settowidth{\lsgbreite}{\,#2\,}%
  \ifdim\lsgbreite>#1\relax\rlap{\raisebox{0.8pt}{\resizebox{#1}{!}{\lsgfarbe\,#2\,}}}%
  \else\rlap{\raisebox{0.8pt}{\lsgfarbe\,#2}}\fi\fi\rule[-0.6ex]{#1}{0.4pt}\xspace}
\newcommand{\luecke}[1][4.5cm]{\lsg[#1]{}}
% Lücke bis Zeilenende: \restzeile / \lsgrest{Text}
\newcommand{\lsgrest}[1]{\leavevmode\unskip\hspace{1.5mm}\ifloesung\rlap{\raisebox{0.8pt}{\lsgfarbe\,#1}}\fi\leaders\hbox{\rule[-0.6ex]{0.5pt}{0.4pt}}\hfill\null\par}
\newcommand{\restzeile}{\lsgrest{}}
% nur in der Lösung sichtbar (Platz bleibt frei): \lsgtext{…}; Alternativen: \lsgoder{Lösung}{Blatt}
\newcommand{\lsgtext}[1]{\ifloesung\leavevmode{\lsgfarbe#1}\fi}
\newcommand{\lsgoder}[2]{\ifloesung\leavevmode#1\else#2\fi}
\newcommand{\lueckentext}{\linespread{1.5}\selectfont\raggedright}
\newcommand{\schreibzeile}{\rule[-2.8mm]{0pt}{8mm}}
% Schritt-Nummern wie im Code und auf der Website: \nr{1} = ①
\newcommand{\nr}[1]{\ding{\numexpr171+#1\relax}}
% Schreiblinien: \linien{n} / \lsglinien{n}{Text} (Text liegt auf den Linien)
\newcommand{\lsglinien}[2]{\par\vspace{0.5mm}\noindent\begin{tikzpicture}
  \path (0,0) (\linewidth-0.4pt,-#1\zeilenabstand-1.5mm);
  \foreach \i in {1,...,#1}{\draw[schreiblinie] (0.2pt,-\i\zeilenabstand) -- (\linewidth-0.2pt,-\i\zeilenabstand);}
  \ifloesung\node[overlay,anchor=north west,inner sep=0pt] at (1mm,-\zeilenabstand+1pt+5mm) {\parbox[t]{\linewidth-2mm}{\lsgfarbe\raggedright\setlength{\baselineskip}{\zeilenabstand}\rule{0pt}{5mm}#2}};\fi
  \end{tikzpicture}\par}
\newcommand{\linien}[1]{\lsglinien{#1}{}}
% Codefeld mit Linien: \linienfeld{n} / \lsgfeld{n}{Code, Zeilen mit \\ getrennt}
\newcommand{\lsgfeld}[2]{\noindent\begin{tikzpicture}
  \draw[line width=0.4pt] (0.2pt,0) rectangle (\linewidth-0.2pt,-#1\zeilenabstand-3mm);
  \foreach \i in {1,...,#1}{\draw[schreiblinie] (2mm,-\i\zeilenabstand) -- (\linewidth-2mm,-\i\zeilenabstand);}
  \ifloesung\node[overlay,anchor=north west,inner sep=0pt] at (2.5mm,-\zeilenabstand+1pt+5mm) {\parbox[t]{\linewidth-5mm}{\lsgfarbe\ttfamily\small\setlength{\baselineskip}{\zeilenabstand}\rule{0pt}{5mm}#2}};\fi
  \end{tikzpicture}}
\newcommand{\linienfeld}[1]{\lsgfeld{#1}{}}
% Karoraster; Lösungszeichnung als TikZ-Code, Ursprung oben links, y nach unten negativ:
% \kariert{n} / \lsgkariert{n}{TikZ},  \karobox{b}{h} / \lsgkarobox{b}{h}{TikZ}
\newcommand{\karoraster}[3]{\begin{tikzpicture}\draw[karo] (0,0) grid[step=\karo] (#1\karo,#2\karo);
  \ifloesung\begin{scope}[overlay,shift={(0,#2\karo)},red!75!black,text=red!75!black,line width=0.8pt]#3\end{scope}\fi\end{tikzpicture}}
\newcommand{\lsgkariert}[2]{\par\vspace{1.5mm}\noindent\pgfmathtruncatemacro{\karoX}{floor(\linewidth/\karo)}%
  \makebox[\linewidth]{\karoraster{\karoX}{#1}{#2}}\par\vspace{1mm}}
\newcommand{\kariert}[1]{\lsgkariert{#1}{}}
\newcommand{\lsgkarobox}[3]{\pgfmathtruncatemacro{\karoX}{round((#1)/\karo)}\pgfmathtruncatemacro{\karoY}{round((#2)/\karo)}\karoraster{\karoX}{\karoY}{#3}}
\newcommand{\karobox}[2]{\lsgkarobox{#1}{#2}{}}
% Lösungs-Klassenkarte (für Zeichnungen im Karo): \lsgkarte{Name}{Attribute}{Methoden}
\newcommand{\lsgkarte}[3]{{\lsgfarbe\setlength{\arrayrulewidth}{0.6pt}\arrayrulecolor{red!75!black}\begin{tabular}{|l|}\hline\textbf{#1}\\\hline #2\\\hline #3\\\hline\end{tabular}}}
% Lösungs-Objektkarte (in TikZ als Knoteninhalt): \node[objekt] {\lsgobjekt{name : Klasse}{attr = wert\\…}};
\newcommand{\lsgobjekt}[2]{{\lsgfarbe\small\arrayrulecolor{red!75!black}\begin{tabular}{@{}l@{}}\textbf{#1}\\\hline #2\end{tabular}}}
\tikzset{objekt/.style={draw,rounded corners=3mm,inner sep=2mm,anchor=north west}}
\newcommand{\klassenkarte}[3]{\begin{tabular}{|l|}\hline\textbf{#1}\\\hline #2\\\hline #3\\\hline\end{tabular}}
\newcommand{\code}[1]{\texttt{#1}}

% ---------- Java-Listings ----------
\lstset{language=Java,basicstyle=\ttfamily\small,keywordstyle=\bfseries,commentstyle=\color{gray!80!black},stringstyle=\color{black},showstringspaces=false,columns=flexible,keepspaces=true,tabsize=3,frame=single,framerule=0.4pt,rulecolor=\color{black},xleftmargin=2mm,framexleftmargin=1mm,aboveskip=2mm,belowskip=2mm,breaklines=true,
 literate={ä}{{\"a}}1 {ö}{{\"o}}1 {ü}{{\"u}}1 {Ä}{{\"A}}1 {Ö}{{\"O}}1 {Ü}{{\"U}}1 {ß}{{\ss}}1 {–}{{--}}1 {„}{{\glqq}}1 {“}{{\grqq}}1}
\lstnewenvironment{java}[1][]{\lstset{#1}}{}
% Lückencode: größerer Zeilenabstand, Lücken mit (*\luecke[3cm]*) im Code
\lstnewenvironment{javaluecke}[1][]{\lstset{basicstyle=\linespread{1.5}\small\ttfamily,escapeinside={(*}{*)},#1}}{}

%  Übersetzen mit XeLaTeX, LuaLaTeX, pdfLaTeX oder Tectonic.
%
%  Blatt:   tectonic ab-5-1-grundlagen.tex   (oder xelatex)
%  Lösung:  tectonic ab-5-1-grundlagen-lsg.tex
%           (Datei enthält nur: \def\mitloesung{}\documentclass[11pt,a4paper]{article}
\input{ab-vorlage}
\newcolumntype{C}{>{\centering\arraybackslash}X}
\newcommand{\kopfzelle}[1]{\cellcolor{blue!8}\textbf{#1}}
% Schlüssel-Symbol
\newcommand{\schluessel}[1][black]{\tikz[baseline=-0.6ex,#1,line width=0.9pt]{\draw (0,0) circle (0.11); \draw (0.11,0) -- (0.5,0); \draw (0.38,0) -- (0.38,-0.1) (0.46,0) -- (0.46,-0.1);}}
\begin{document}
\abtitel{5.1}{Grundlagen der Verschlüsselung}

\begin{aufgabe}[Die Skytale]{1}
\begin{minipage}[t]{0.58\linewidth}
\vspace{0pt}
Die Spartaner wickelten vor 2500 Jahren einen Lederstreifen um einen Stab (die \textbf{Skytale}) und schrieben ihre Nachricht der Länge nach darauf. Abgewickelt sieht der Streifen wie Buchstabensalat aus. Nur wer einen gleich dicken Stab hat, kann die Nachricht lesen.\\[1mm]
Abgefangener Streifen: \quad\texttt{\large TKLHEIEIECSXAAX}\\[1mm]
Um den Stab passen genau \textbf{3} Buchstaben. Schreibe den Streifen \textbf{spaltenweise} (von oben nach unten) in das Raster und lies \textbf{zeilenweise}.\par
{\lueckentext Klartext: \lsg[6cm]{THE CAKE IS A LIE}\par}
\end{minipage}\hfill
\begin{minipage}[t]{0.38\linewidth}
\vspace{2mm}\centering
\begin{tikzpicture}[x=7mm,y=7mm]
  \draw[gray!70] (0,0) grid (5,3);
  \draw[line width=0.7pt] (0,0) rectangle (5,3);
  \node[font=\small,anchor=east] at (-0.1,2.5) {Zeile 1};
  \node[font=\small,anchor=east] at (-0.1,1.5) {Zeile 2};
  \node[font=\small,anchor=east] at (-0.1,0.5) {Zeile 3};
  \node[font=\ttfamily] at (0.5,2.5) {T}; \node[font=\ttfamily] at (0.5,1.5) {K}; \node[font=\ttfamily] at (0.5,0.5) {L};
  \draw[-{Stealth},blue!60!black,thick] (0.5,3.75) -- (0.5,3.15);
  \ifloesung
  \foreach \z [count=\r] in {{T,H,E,C,A},{K,E,I,S,A},{L,I,E,X,X}} {
    \foreach \b [count=\s] in \z {\ifnum\s>1 \node[font=\ttfamily,text=red!75!black] at (\s-0.5,3.5-\r) {\b};\fi}
  }
  \fi
\end{tikzpicture}\\[1mm]
{\small (X = Füllzeichen)}
\end{minipage}
\end{aufgabe}

\begin{minipage}[t]{0.63\linewidth}
\vspace{0pt}
Helga ist bei einer Recherche zu \textbf{Verschlüsselungen} auf einen Forumseintrag gestoßen. Die Antwort wurde mit dem \textbf{ROT13-Verfahren} verschlüsselt: Jeder Buchstabe des Alphabets wird um 13 Stellen weitergeschoben (nach Z geht es bei A weiter).
\end{minipage}\hfill
\begin{minipage}[t]{0.33\linewidth}
\vspace{0pt}
\begin{tcolorbox}[colback=yellow!15,colframe=orange!60,boxrule=0.5pt,arc=1mm,left=1.5mm,right=1.5mm,top=1mm,bottom=1mm]
\small\itshape Wem wird die Entschlüsselung der Vigenère-Chiffre im 19.\,Jahrhundert zugeschrieben?\\[0.5mm]
\upshape Antwort: \texttt{\bfseries PUNEYRF ONOONTR}
\end{tcolorbox}
\end{minipage}

\begin{aufgabe}[ROT13]{2}
\textbf{a)} Ergänze die Tabelle: Unter jedem Buchstaben steht der Buchstabe, der 13 Stellen weiter liegt. Entschlüssle dann Helgas Antwort.\par\vspace{1.5mm}
{\setlength{\tabcolsep}{0pt}\renewcommand{\arraystretch}{1.25}\small
\begin{tabularx}{\linewidth}{|*{26}{C|}}\hline
A&B&C&D&E&F&G&H&I&J&K&L&M&N&O&P&Q&R&S&T&U&V&W&X&Y&Z\\\hline
\schreibzeile\lsgtext{N}&\lsgtext{O}&\lsgtext{P}&\lsgtext{Q}&\lsgtext{R}&\lsgtext{S}&\lsgtext{T}&\lsgtext{U}&\lsgtext{V}&\lsgtext{W}&\lsgtext{X}&\lsgtext{Y}&\lsgtext{Z}&\lsgtext{A}&\lsgtext{B}&\lsgtext{C}&\lsgtext{D}&\lsgtext{E}&\lsgtext{F}&\lsgtext{G}&\lsgtext{H}&\lsgtext{I}&\lsgtext{J}&\lsgtext{K}&\lsgtext{L}&\lsgtext{M}\\\hline
\end{tabularx}}\par\vspace{1mm}
{\lueckentext Klartext: \lsg[6cm]{CHARLES BABBAGE}\\
\textbf{b)} Öffne das Projekt \texttt{11-5-1 Verschluesselung} (Website) und erklärt euch gegenseitig, wie die Methoden \texttt{rotiere} und \texttt{rot13} arbeiten. Ergänze:\\
Der Computer rechnet mit dem ASCII-Code: \texttt{'A'} hat den Code 65, \texttt{'Z'} den Code \lsg[1.5cm]{90}. In \nr{2} wird die Verschiebung zum Code addiert. Ist das Ergebnis größer als 90, werden in \nr{3} \lsg[1.5cm]{26} abgezogen. In \texttt{rot13} geht eine \lsg[2.5cm]{for-Schleife} den Text Zeichen für Zeichen durch, verschiebt jedes Zeichen um \lsg[1.5cm]{13} und hängt es hinten an den Geheimtext an.\\
\textbf{c)} Verschlüssle den Geheimtext noch einmal mit ROT13. Warum entsteht wieder der Klartext?\par}
\lsglinien{1}{13 + 13 = 26 Stellen: einmal ganz um das Alphabet herum. Ver- und Entschlüsseln sind gleich.}
\end{aufgabe}

\newpage
\textbf{Grundprinzip der Verschlüsselung}\par\vspace{1mm}
\begin{minipage}[c]{0.5\linewidth}
\centering
\begin{tikzpicture}[font=\small,
  txt/.style={draw,ellipse,minimum width=16mm,minimum height=7mm,inner sep=1pt}]
  \node[txt,fill=green!15] (k1) at (0,0) {Klartext};
  \node[txt,fill=orange!20] (g) at (3.05,0) {Geheimtext};
  \node[txt,fill=green!15] (k2) at (6.1,0) {Klartext};
  \draw[-{Stealth},thick] (k1.north east) to[bend left=40] node[above,align=center,font=\scriptsize] {\scalebox{1.6}{\schluessel}\\Verschlüsselung} (g.north west);
  \draw[-{Stealth},thick] (g.north east) to[bend left=40] node[above,align=center,font=\scriptsize] {\scalebox{1.6}{\schluessel}\\Entschlüsselung} (k2.north west);
  \draw[red!70!black] (-0.8,-0.6) -- (-0.8,-0.75) -- (2.1,-0.75) -- (2.1,-0.6); \node[red!70!black,font=\scriptsize] at (0.65,-1) {Sender};
  \draw[violet] (4.0,-0.6) -- (4.0,-0.75) -- (6.9,-0.75) -- (6.9,-0.6); \node[violet,font=\scriptsize] at (5.45,-1) {Empfänger};
\end{tikzpicture}
\end{minipage}\hfill
\begin{minipage}[c]{0.47\linewidth}
{\lueckentext Im Gegensatz zur Codierung (Kapitel 4: Regeln öffentlich) wird bei der \lsg[3.8cm]{Verschlüsselung} eine Nachricht (\textbf{Klartext}) in unlesbare Daten (\textbf{Geheimtext}) umgewandelt, sodass sie vor Unbefugten geschützt ist. Nur wer zusätzlich den \lsg[3cm]{Schlüssel} kennt, kann den Geheimtext wieder entschlüsseln.\par}
\end{minipage}

\medskip
\textbf{Sicherheitsziele einer Verschlüsselung}\par
\begin{tabularx}{\linewidth}{@{}l X@{}}
\textbf{Vertraulichkeit:} & Eine Nachricht kann nur von Befugten gelesen werden.\\
\textbf{Integrität:} & Eine Nachricht kann nicht unbemerkt verändert werden.\\
\textbf{Authentizität:} & Eine Nachricht kann nicht unter falschem Namen verschickt werden.\\
\textbf{Verbindlichkeit:} & Der Absender kann später nicht bestreiten, die Nachricht verfasst zu haben.\\
\end{tabularx}

\begin{aufgabe}[Welches Ziel ist in Gefahr?]{3}
{\lueckentext
\nr{1} Im Schul-WLAN liest jemand deine Chatnachrichten mit: \lsgrest{Vertraulichkeit}
\nr{2} Bei einer Online-Überweisung werden unterwegs aus 10\,€ plötzlich 100\,€: \lsgrest{Integrität}
\nr{3} Eine E-Mail sieht aus, als käme sie von deiner Bank, stammt aber von Betrügern: \lsgrest{Authentizität}
\nr{4} Jemand bestellt online und behauptet später, er habe nie bestellt: \lsgrest{Verbindlichkeit}
\par}
\end{aufgabe}

\textbf{Ein Schlüssel oder ein Schlüsselpaar?}\par
{\lueckentext Man unterscheidet vor allem zwei Arten von Verfahren:\\
\lsg[3.8cm]{Symmetrische} Verschlüsselung: Derselbe Schlüssel wird zum Ver- und Entschlüsseln benutzt.\\
\lsg[3.8cm]{Asymmetrische} Verschlüsselung: Es gibt ein \textbf{Schlüsselpaar}. Mit dem einen Schlüssel kann man nur verschlüsseln, mit dem anderen nur entschlüsseln.\par}

\begin{aufgabe}[The Mystery of Cryptocastle]{4}
Öffne das Spiel (Link auf der Website), gehe zu Abschnitt \textbf{1} und löse die Aufgaben bis zu der Truhe, in der sich der Code \texttt{\bfseries GOVTGYZTTMZKCYE} befindet. Den Code knackst du in Station 5.2.
\end{aufgabe}

\vfill
\web{Station 5.1 Grundlagen: Skytale, Chiffrierscheibe, Code-Animation zu rot13}
\end{document}
;
%            füllt alle Lücken rot aus; siehe \lsg, \lsgrest, \lsglinien,
%            \lsgfeld, \lsgkariert, \lsgkarobox, \lsgtext, \lsgoder)
%  Lösungseinträge belegen nie mehr Platz als die Lücke im Blatt.
% =====================================================================
\usepackage[a4paper,left=18mm,right=18mm,top=26mm,bottom=17mm,headheight=15mm,headsep=4mm,footskip=8mm]{geometry}
\usepackage{iftex}
\ifPDFTeX
  \usepackage[T1]{fontenc}
  \usepackage[utf8]{inputenc}
  \usepackage{helvet}
  \usepackage[scaled=0.86]{beramono}
  \renewcommand{\familydefault}{\sfdefault}
\else
  \usepackage{fontspec}
  \setmainfont{texgyreheros}[Extension=.otf,UprightFont=*-regular,BoldFont=*-bold,ItalicFont=*-italic,BoldItalicFont=*-bolditalic]
  \setmonofont{DejaVuSansMono}[Extension=.ttf,UprightFont=*,BoldFont=*-Bold,ItalicFont=*-Oblique,BoldItalicFont=*-BoldOblique,Scale=0.86]
\fi
\usepackage[ngerman,shorthands=off]{babel}
\usepackage{xcolor,graphicx,tabularx,array,enumitem,fancyhdr,tikz,listings,multicol,calc,etoolbox,ragged2e,colortbl,xspace,pifont}
\usepackage[most]{tcolorbox}
\usetikzlibrary{arrows.meta,positioning,calc,decorations.pathreplacing,shapes.geometric,shapes.symbols,fit}
\usepackage[hidelinks]{hyperref}
\setlength{\parindent}{0pt}
\setlength{\parskip}{3pt}
\setlist{nosep,leftmargin=*}

% ---------- Lösungsschalter ----------
\newif\ifloesung
\ifdefined\mitloesung\loesungtrue\fi
\newcommand{\lsgfarbe}{\color{red!75!black}}

% ---------- Kopf- und Fußzeile (einheitliche Form) ----------
\newcommand{\lehrkraft}{Hau}
\newcommand{\themenbereich}{5. Verschlüsselung}
\newcommand{\abnummer}{}
\pagestyle{fancy}\fancyhf{}
\renewcommand{\headrulewidth}{0.4pt}
\fancyhead[L]{\small Klasse 11\\Datum:}
\fancyhead[C]{\small Themenbereich:\\\themenbereich}
\fancyhead[R]{\small \lehrkraft\ifloesung\\{\lsgfarbe\bfseries Lösung}\fi}
\fancyfoot[C]{\scriptsize edu-mrh.de\,·\,Informatik 11\,·\,Blatt \abnummer\ifloesung\ (Lösung)\fi\,·\,Seite \thepage}
\newcommand{\abtitel}[2]{\renewcommand{\abnummer}{#1}\begin{center}\Large\bfseries #1\quad\underline{#2}\end{center}\vspace{-1mm}}

% ---------- Boxen ----------
\newenvironment{aufgabe}[2][]{\begin{tcolorbox}[enhanced,breakable,colback=white,colframe=black,boxrule=0.7pt,arc=3mm,left=2mm,right=2mm,top=1.5mm,bottom=1.5mm,before skip=2.5mm,after skip=2.5mm]\textbf{Aufgabe #2}\ifstrempty{#1}{}{ \textit{#1}}\par\vspace{0.6mm}}{\end{tcolorbox}}
\newtcolorbox{merke}[1][Merke]{enhanced,breakable,colback=green!5,colframe=green!40!black,boxrule=0.8pt,arc=2mm,left=2mm,right=2mm,top=1.5mm,bottom=1.5mm,before skip=2.5mm,after skip=2.5mm,title={#1},fonttitle=\bfseries,coltitle=white}
\newtcolorbox{hinweis}{enhanced,breakable,colback=yellow!12,colframe=orange!70!black,boxrule=0.6pt,arc=2mm,left=2mm,right=2mm,top=1mm,bottom=1mm,before skip=2mm,after skip=2mm}
\newcommand{\web}[1]{{\small\color{gray}\textit{Website: #1}}}
% Kasten für optionale Aufgaben
\newtcolorbox{optional}[1][Optional]{enhanced,breakable,colback=blue!3,colframe=blue!45!black,boxrule=0.6pt,arc=2mm,left=2mm,right=2mm,top=1.5mm,bottom=1.5mm,before skip=2.5mm,after skip=2.5mm,title={#1},fonttitle=\bfseries\small,coltitle=white}
% Zahl mit Basis: \zs{1011}{2}
\newcommand{\zs}[2]{#1\textsubscript{#2}}

% ---------- Lücken, Linien, Karos ----------
% Einheitliche Maße für alle Blätter der Sequenz:
%  \linien{n}       n Schreiblinien, Abstand 8 mm, erste Linie 8 mm unter dem Text
%  \kariert{n}      Karoraster 5 mm, n Kästchen hoch, nur ganze Kästchen, mittig
%  \karobox{b}{h}   Karoraster b x h, auf ganze Kästchen gerundet
%  \luecke[l]       Lücke im Fließtext (Standard 4,5 cm, Fachbegriffe mind. 3,5 cm)
%  (jeweils mit Lösungsvariante \lsg…, siehe Kopf der Datei)
%  \restzeile       Lücke bis zum Zeilenende
%  \lueckentext     am Anfang eines Absatzes/Kastens mit Lücken: größerer
%                   Zeilenabstand und Flattersatz, damit man in die Lücken schreiben kann
%  \linienfeld{n}   umrahmtes Feld mit n Schreiblinien (z. B. für Java-Code)
%  \schreibzeile    Strut für Tabellenzellen, in die geschrieben wird (8 mm hoch)
\newlength{\zeilenabstand}\setlength{\zeilenabstand}{8mm}
\newlength{\karo}\setlength{\karo}{5mm}
\tikzset{karo/.style={draw=gray!50,line width=0.3pt},schreiblinie/.style={draw=black!65,line width=0.35pt}}
\newlength{\lsgbreite}
% Lücke: \luecke[Breite] / mit Lösung \lsg[Breite]{Text} (zu langer Text wird verkleinert)
\newcommand{\lsg}[2][4.5cm]{\leavevmode\ifloesung\settowidth{\lsgbreite}{\,#2\,}%
  \ifdim\lsgbreite>#1\relax\rlap{\raisebox{0.8pt}{\resizebox{#1}{!}{\lsgfarbe\,#2\,}}}%
  \else\rlap{\raisebox{0.8pt}{\lsgfarbe\,#2}}\fi\fi\rule[-0.6ex]{#1}{0.4pt}\xspace}
\newcommand{\luecke}[1][4.5cm]{\lsg[#1]{}}
% Lücke bis Zeilenende: \restzeile / \lsgrest{Text}
\newcommand{\lsgrest}[1]{\leavevmode\unskip\hspace{1.5mm}\ifloesung\rlap{\raisebox{0.8pt}{\lsgfarbe\,#1}}\fi\leaders\hbox{\rule[-0.6ex]{0.5pt}{0.4pt}}\hfill\null\par}
\newcommand{\restzeile}{\lsgrest{}}
% nur in der Lösung sichtbar (Platz bleibt frei): \lsgtext{…}; Alternativen: \lsgoder{Lösung}{Blatt}
\newcommand{\lsgtext}[1]{\ifloesung\leavevmode{\lsgfarbe#1}\fi}
\newcommand{\lsgoder}[2]{\ifloesung\leavevmode#1\else#2\fi}
\newcommand{\lueckentext}{\linespread{1.5}\selectfont\raggedright}
\newcommand{\schreibzeile}{\rule[-2.8mm]{0pt}{8mm}}
% Schritt-Nummern wie im Code und auf der Website: \nr{1} = ①
\newcommand{\nr}[1]{\ding{\numexpr171+#1\relax}}
% Schreiblinien: \linien{n} / \lsglinien{n}{Text} (Text liegt auf den Linien)
\newcommand{\lsglinien}[2]{\par\vspace{0.5mm}\noindent\begin{tikzpicture}
  \path (0,0) (\linewidth-0.4pt,-#1\zeilenabstand-1.5mm);
  \foreach \i in {1,...,#1}{\draw[schreiblinie] (0.2pt,-\i\zeilenabstand) -- (\linewidth-0.2pt,-\i\zeilenabstand);}
  \ifloesung\node[overlay,anchor=north west,inner sep=0pt] at (1mm,-\zeilenabstand+1pt+5mm) {\parbox[t]{\linewidth-2mm}{\lsgfarbe\raggedright\setlength{\baselineskip}{\zeilenabstand}\rule{0pt}{5mm}#2}};\fi
  \end{tikzpicture}\par}
\newcommand{\linien}[1]{\lsglinien{#1}{}}
% Codefeld mit Linien: \linienfeld{n} / \lsgfeld{n}{Code, Zeilen mit \\ getrennt}
\newcommand{\lsgfeld}[2]{\noindent\begin{tikzpicture}
  \draw[line width=0.4pt] (0.2pt,0) rectangle (\linewidth-0.2pt,-#1\zeilenabstand-3mm);
  \foreach \i in {1,...,#1}{\draw[schreiblinie] (2mm,-\i\zeilenabstand) -- (\linewidth-2mm,-\i\zeilenabstand);}
  \ifloesung\node[overlay,anchor=north west,inner sep=0pt] at (2.5mm,-\zeilenabstand+1pt+5mm) {\parbox[t]{\linewidth-5mm}{\lsgfarbe\ttfamily\small\setlength{\baselineskip}{\zeilenabstand}\rule{0pt}{5mm}#2}};\fi
  \end{tikzpicture}}
\newcommand{\linienfeld}[1]{\lsgfeld{#1}{}}
% Karoraster; Lösungszeichnung als TikZ-Code, Ursprung oben links, y nach unten negativ:
% \kariert{n} / \lsgkariert{n}{TikZ},  \karobox{b}{h} / \lsgkarobox{b}{h}{TikZ}
\newcommand{\karoraster}[3]{\begin{tikzpicture}\draw[karo] (0,0) grid[step=\karo] (#1\karo,#2\karo);
  \ifloesung\begin{scope}[overlay,shift={(0,#2\karo)},red!75!black,text=red!75!black,line width=0.8pt]#3\end{scope}\fi\end{tikzpicture}}
\newcommand{\lsgkariert}[2]{\par\vspace{1.5mm}\noindent\pgfmathtruncatemacro{\karoX}{floor(\linewidth/\karo)}%
  \makebox[\linewidth]{\karoraster{\karoX}{#1}{#2}}\par\vspace{1mm}}
\newcommand{\kariert}[1]{\lsgkariert{#1}{}}
\newcommand{\lsgkarobox}[3]{\pgfmathtruncatemacro{\karoX}{round((#1)/\karo)}\pgfmathtruncatemacro{\karoY}{round((#2)/\karo)}\karoraster{\karoX}{\karoY}{#3}}
\newcommand{\karobox}[2]{\lsgkarobox{#1}{#2}{}}
% Lösungs-Klassenkarte (für Zeichnungen im Karo): \lsgkarte{Name}{Attribute}{Methoden}
\newcommand{\lsgkarte}[3]{{\lsgfarbe\setlength{\arrayrulewidth}{0.6pt}\arrayrulecolor{red!75!black}\begin{tabular}{|l|}\hline\textbf{#1}\\\hline #2\\\hline #3\\\hline\end{tabular}}}
% Lösungs-Objektkarte (in TikZ als Knoteninhalt): \node[objekt] {\lsgobjekt{name : Klasse}{attr = wert\\…}};
\newcommand{\lsgobjekt}[2]{{\lsgfarbe\small\arrayrulecolor{red!75!black}\begin{tabular}{@{}l@{}}\textbf{#1}\\\hline #2\end{tabular}}}
\tikzset{objekt/.style={draw,rounded corners=3mm,inner sep=2mm,anchor=north west}}
\newcommand{\klassenkarte}[3]{\begin{tabular}{|l|}\hline\textbf{#1}\\\hline #2\\\hline #3\\\hline\end{tabular}}
\newcommand{\code}[1]{\texttt{#1}}

% ---------- Java-Listings ----------
\lstset{language=Java,basicstyle=\ttfamily\small,keywordstyle=\bfseries,commentstyle=\color{gray!80!black},stringstyle=\color{black},showstringspaces=false,columns=flexible,keepspaces=true,tabsize=3,frame=single,framerule=0.4pt,rulecolor=\color{black},xleftmargin=2mm,framexleftmargin=1mm,aboveskip=2mm,belowskip=2mm,breaklines=true,
 literate={ä}{{\"a}}1 {ö}{{\"o}}1 {ü}{{\"u}}1 {Ä}{{\"A}}1 {Ö}{{\"O}}1 {Ü}{{\"U}}1 {ß}{{\ss}}1 {–}{{--}}1 {„}{{\glqq}}1 {“}{{\grqq}}1}
\lstnewenvironment{java}[1][]{\lstset{#1}}{}
% Lückencode: größerer Zeilenabstand, Lücken mit (*\luecke[3cm]*) im Code
\lstnewenvironment{javaluecke}[1][]{\lstset{basicstyle=\linespread{1.5}\small\ttfamily,escapeinside={(*}{*)},#1}}{}

%  Übersetzen mit XeLaTeX, LuaLaTeX, pdfLaTeX oder Tectonic.
%
%  Blatt:   tectonic ab-5-1-grundlagen.tex   (oder xelatex)
%  Lösung:  tectonic ab-5-1-grundlagen-lsg.tex
%           (Datei enthält nur: \def\mitloesung{}\documentclass[11pt,a4paper]{article}
% =====================================================================
%  Gemeinsame Vorlage der Arbeitsblätter
%  "Verschlüsselung" – Informatik 11 (NTG)
%  (Antwortplatz-Makros wie informatik/10/java/arbeitsblaetter/ab-vorlage.tex)
%  Einbinden in jedem Blatt direkt nach \documentclass: \input{ab-vorlage}
%  Übersetzen mit XeLaTeX, LuaLaTeX, pdfLaTeX oder Tectonic.
%
%  Blatt:   tectonic ab-5-1-grundlagen.tex   (oder xelatex)
%  Lösung:  tectonic ab-5-1-grundlagen-lsg.tex
%           (Datei enthält nur: \def\mitloesung{}\input{ab-5-1-grundlagen};
%            füllt alle Lücken rot aus; siehe \lsg, \lsgrest, \lsglinien,
%            \lsgfeld, \lsgkariert, \lsgkarobox, \lsgtext, \lsgoder)
%  Lösungseinträge belegen nie mehr Platz als die Lücke im Blatt.
% =====================================================================
\usepackage[a4paper,left=18mm,right=18mm,top=26mm,bottom=17mm,headheight=15mm,headsep=4mm,footskip=8mm]{geometry}
\usepackage{iftex}
\ifPDFTeX
  \usepackage[T1]{fontenc}
  \usepackage[utf8]{inputenc}
  \usepackage{helvet}
  \usepackage[scaled=0.86]{beramono}
  \renewcommand{\familydefault}{\sfdefault}
\else
  \usepackage{fontspec}
  \setmainfont{texgyreheros}[Extension=.otf,UprightFont=*-regular,BoldFont=*-bold,ItalicFont=*-italic,BoldItalicFont=*-bolditalic]
  \setmonofont{DejaVuSansMono}[Extension=.ttf,UprightFont=*,BoldFont=*-Bold,ItalicFont=*-Oblique,BoldItalicFont=*-BoldOblique,Scale=0.86]
\fi
\usepackage[ngerman,shorthands=off]{babel}
\usepackage{xcolor,graphicx,tabularx,array,enumitem,fancyhdr,tikz,listings,multicol,calc,etoolbox,ragged2e,colortbl,xspace,pifont}
\usepackage[most]{tcolorbox}
\usetikzlibrary{arrows.meta,positioning,calc,decorations.pathreplacing,shapes.geometric,shapes.symbols,fit}
\usepackage[hidelinks]{hyperref}
\setlength{\parindent}{0pt}
\setlength{\parskip}{3pt}
\setlist{nosep,leftmargin=*}

% ---------- Lösungsschalter ----------
\newif\ifloesung
\ifdefined\mitloesung\loesungtrue\fi
\newcommand{\lsgfarbe}{\color{red!75!black}}

% ---------- Kopf- und Fußzeile (einheitliche Form) ----------
\newcommand{\lehrkraft}{Hau}
\newcommand{\themenbereich}{5. Verschlüsselung}
\newcommand{\abnummer}{}
\pagestyle{fancy}\fancyhf{}
\renewcommand{\headrulewidth}{0.4pt}
\fancyhead[L]{\small Klasse 11\\Datum:}
\fancyhead[C]{\small Themenbereich:\\\themenbereich}
\fancyhead[R]{\small \lehrkraft\ifloesung\\{\lsgfarbe\bfseries Lösung}\fi}
\fancyfoot[C]{\scriptsize edu-mrh.de\,·\,Informatik 11\,·\,Blatt \abnummer\ifloesung\ (Lösung)\fi\,·\,Seite \thepage}
\newcommand{\abtitel}[2]{\renewcommand{\abnummer}{#1}\begin{center}\Large\bfseries #1\quad\underline{#2}\end{center}\vspace{-1mm}}

% ---------- Boxen ----------
\newenvironment{aufgabe}[2][]{\begin{tcolorbox}[enhanced,breakable,colback=white,colframe=black,boxrule=0.7pt,arc=3mm,left=2mm,right=2mm,top=1.5mm,bottom=1.5mm,before skip=2.5mm,after skip=2.5mm]\textbf{Aufgabe #2}\ifstrempty{#1}{}{ \textit{#1}}\par\vspace{0.6mm}}{\end{tcolorbox}}
\newtcolorbox{merke}[1][Merke]{enhanced,breakable,colback=green!5,colframe=green!40!black,boxrule=0.8pt,arc=2mm,left=2mm,right=2mm,top=1.5mm,bottom=1.5mm,before skip=2.5mm,after skip=2.5mm,title={#1},fonttitle=\bfseries,coltitle=white}
\newtcolorbox{hinweis}{enhanced,breakable,colback=yellow!12,colframe=orange!70!black,boxrule=0.6pt,arc=2mm,left=2mm,right=2mm,top=1mm,bottom=1mm,before skip=2mm,after skip=2mm}
\newcommand{\web}[1]{{\small\color{gray}\textit{Website: #1}}}
% Kasten für optionale Aufgaben
\newtcolorbox{optional}[1][Optional]{enhanced,breakable,colback=blue!3,colframe=blue!45!black,boxrule=0.6pt,arc=2mm,left=2mm,right=2mm,top=1.5mm,bottom=1.5mm,before skip=2.5mm,after skip=2.5mm,title={#1},fonttitle=\bfseries\small,coltitle=white}
% Zahl mit Basis: \zs{1011}{2}
\newcommand{\zs}[2]{#1\textsubscript{#2}}

% ---------- Lücken, Linien, Karos ----------
% Einheitliche Maße für alle Blätter der Sequenz:
%  \linien{n}       n Schreiblinien, Abstand 8 mm, erste Linie 8 mm unter dem Text
%  \kariert{n}      Karoraster 5 mm, n Kästchen hoch, nur ganze Kästchen, mittig
%  \karobox{b}{h}   Karoraster b x h, auf ganze Kästchen gerundet
%  \luecke[l]       Lücke im Fließtext (Standard 4,5 cm, Fachbegriffe mind. 3,5 cm)
%  (jeweils mit Lösungsvariante \lsg…, siehe Kopf der Datei)
%  \restzeile       Lücke bis zum Zeilenende
%  \lueckentext     am Anfang eines Absatzes/Kastens mit Lücken: größerer
%                   Zeilenabstand und Flattersatz, damit man in die Lücken schreiben kann
%  \linienfeld{n}   umrahmtes Feld mit n Schreiblinien (z. B. für Java-Code)
%  \schreibzeile    Strut für Tabellenzellen, in die geschrieben wird (8 mm hoch)
\newlength{\zeilenabstand}\setlength{\zeilenabstand}{8mm}
\newlength{\karo}\setlength{\karo}{5mm}
\tikzset{karo/.style={draw=gray!50,line width=0.3pt},schreiblinie/.style={draw=black!65,line width=0.35pt}}
\newlength{\lsgbreite}
% Lücke: \luecke[Breite] / mit Lösung \lsg[Breite]{Text} (zu langer Text wird verkleinert)
\newcommand{\lsg}[2][4.5cm]{\leavevmode\ifloesung\settowidth{\lsgbreite}{\,#2\,}%
  \ifdim\lsgbreite>#1\relax\rlap{\raisebox{0.8pt}{\resizebox{#1}{!}{\lsgfarbe\,#2\,}}}%
  \else\rlap{\raisebox{0.8pt}{\lsgfarbe\,#2}}\fi\fi\rule[-0.6ex]{#1}{0.4pt}\xspace}
\newcommand{\luecke}[1][4.5cm]{\lsg[#1]{}}
% Lücke bis Zeilenende: \restzeile / \lsgrest{Text}
\newcommand{\lsgrest}[1]{\leavevmode\unskip\hspace{1.5mm}\ifloesung\rlap{\raisebox{0.8pt}{\lsgfarbe\,#1}}\fi\leaders\hbox{\rule[-0.6ex]{0.5pt}{0.4pt}}\hfill\null\par}
\newcommand{\restzeile}{\lsgrest{}}
% nur in der Lösung sichtbar (Platz bleibt frei): \lsgtext{…}; Alternativen: \lsgoder{Lösung}{Blatt}
\newcommand{\lsgtext}[1]{\ifloesung\leavevmode{\lsgfarbe#1}\fi}
\newcommand{\lsgoder}[2]{\ifloesung\leavevmode#1\else#2\fi}
\newcommand{\lueckentext}{\linespread{1.5}\selectfont\raggedright}
\newcommand{\schreibzeile}{\rule[-2.8mm]{0pt}{8mm}}
% Schritt-Nummern wie im Code und auf der Website: \nr{1} = ①
\newcommand{\nr}[1]{\ding{\numexpr171+#1\relax}}
% Schreiblinien: \linien{n} / \lsglinien{n}{Text} (Text liegt auf den Linien)
\newcommand{\lsglinien}[2]{\par\vspace{0.5mm}\noindent\begin{tikzpicture}
  \path (0,0) (\linewidth-0.4pt,-#1\zeilenabstand-1.5mm);
  \foreach \i in {1,...,#1}{\draw[schreiblinie] (0.2pt,-\i\zeilenabstand) -- (\linewidth-0.2pt,-\i\zeilenabstand);}
  \ifloesung\node[overlay,anchor=north west,inner sep=0pt] at (1mm,-\zeilenabstand+1pt+5mm) {\parbox[t]{\linewidth-2mm}{\lsgfarbe\raggedright\setlength{\baselineskip}{\zeilenabstand}\rule{0pt}{5mm}#2}};\fi
  \end{tikzpicture}\par}
\newcommand{\linien}[1]{\lsglinien{#1}{}}
% Codefeld mit Linien: \linienfeld{n} / \lsgfeld{n}{Code, Zeilen mit \\ getrennt}
\newcommand{\lsgfeld}[2]{\noindent\begin{tikzpicture}
  \draw[line width=0.4pt] (0.2pt,0) rectangle (\linewidth-0.2pt,-#1\zeilenabstand-3mm);
  \foreach \i in {1,...,#1}{\draw[schreiblinie] (2mm,-\i\zeilenabstand) -- (\linewidth-2mm,-\i\zeilenabstand);}
  \ifloesung\node[overlay,anchor=north west,inner sep=0pt] at (2.5mm,-\zeilenabstand+1pt+5mm) {\parbox[t]{\linewidth-5mm}{\lsgfarbe\ttfamily\small\setlength{\baselineskip}{\zeilenabstand}\rule{0pt}{5mm}#2}};\fi
  \end{tikzpicture}}
\newcommand{\linienfeld}[1]{\lsgfeld{#1}{}}
% Karoraster; Lösungszeichnung als TikZ-Code, Ursprung oben links, y nach unten negativ:
% \kariert{n} / \lsgkariert{n}{TikZ},  \karobox{b}{h} / \lsgkarobox{b}{h}{TikZ}
\newcommand{\karoraster}[3]{\begin{tikzpicture}\draw[karo] (0,0) grid[step=\karo] (#1\karo,#2\karo);
  \ifloesung\begin{scope}[overlay,shift={(0,#2\karo)},red!75!black,text=red!75!black,line width=0.8pt]#3\end{scope}\fi\end{tikzpicture}}
\newcommand{\lsgkariert}[2]{\par\vspace{1.5mm}\noindent\pgfmathtruncatemacro{\karoX}{floor(\linewidth/\karo)}%
  \makebox[\linewidth]{\karoraster{\karoX}{#1}{#2}}\par\vspace{1mm}}
\newcommand{\kariert}[1]{\lsgkariert{#1}{}}
\newcommand{\lsgkarobox}[3]{\pgfmathtruncatemacro{\karoX}{round((#1)/\karo)}\pgfmathtruncatemacro{\karoY}{round((#2)/\karo)}\karoraster{\karoX}{\karoY}{#3}}
\newcommand{\karobox}[2]{\lsgkarobox{#1}{#2}{}}
% Lösungs-Klassenkarte (für Zeichnungen im Karo): \lsgkarte{Name}{Attribute}{Methoden}
\newcommand{\lsgkarte}[3]{{\lsgfarbe\setlength{\arrayrulewidth}{0.6pt}\arrayrulecolor{red!75!black}\begin{tabular}{|l|}\hline\textbf{#1}\\\hline #2\\\hline #3\\\hline\end{tabular}}}
% Lösungs-Objektkarte (in TikZ als Knoteninhalt): \node[objekt] {\lsgobjekt{name : Klasse}{attr = wert\\…}};
\newcommand{\lsgobjekt}[2]{{\lsgfarbe\small\arrayrulecolor{red!75!black}\begin{tabular}{@{}l@{}}\textbf{#1}\\\hline #2\end{tabular}}}
\tikzset{objekt/.style={draw,rounded corners=3mm,inner sep=2mm,anchor=north west}}
\newcommand{\klassenkarte}[3]{\begin{tabular}{|l|}\hline\textbf{#1}\\\hline #2\\\hline #3\\\hline\end{tabular}}
\newcommand{\code}[1]{\texttt{#1}}

% ---------- Java-Listings ----------
\lstset{language=Java,basicstyle=\ttfamily\small,keywordstyle=\bfseries,commentstyle=\color{gray!80!black},stringstyle=\color{black},showstringspaces=false,columns=flexible,keepspaces=true,tabsize=3,frame=single,framerule=0.4pt,rulecolor=\color{black},xleftmargin=2mm,framexleftmargin=1mm,aboveskip=2mm,belowskip=2mm,breaklines=true,
 literate={ä}{{\"a}}1 {ö}{{\"o}}1 {ü}{{\"u}}1 {Ä}{{\"A}}1 {Ö}{{\"O}}1 {Ü}{{\"U}}1 {ß}{{\ss}}1 {–}{{--}}1 {„}{{\glqq}}1 {“}{{\grqq}}1}
\lstnewenvironment{java}[1][]{\lstset{#1}}{}
% Lückencode: größerer Zeilenabstand, Lücken mit (*\luecke[3cm]*) im Code
\lstnewenvironment{javaluecke}[1][]{\lstset{basicstyle=\linespread{1.5}\small\ttfamily,escapeinside={(*}{*)},#1}}{}

\newcolumntype{C}{>{\centering\arraybackslash}X}
\newcommand{\kopfzelle}[1]{\cellcolor{blue!8}\textbf{#1}}
% Schlüssel-Symbol
\newcommand{\schluessel}[1][black]{\tikz[baseline=-0.6ex,#1,line width=0.9pt]{\draw (0,0) circle (0.11); \draw (0.11,0) -- (0.5,0); \draw (0.38,0) -- (0.38,-0.1) (0.46,0) -- (0.46,-0.1);}}
\begin{document}
\abtitel{5.1}{Grundlagen der Verschlüsselung}

\begin{aufgabe}[Die Skytale]{1}
\begin{minipage}[t]{0.58\linewidth}
\vspace{0pt}
Die Spartaner wickelten vor 2500 Jahren einen Lederstreifen um einen Stab (die \textbf{Skytale}) und schrieben ihre Nachricht der Länge nach darauf. Abgewickelt sieht der Streifen wie Buchstabensalat aus. Nur wer einen gleich dicken Stab hat, kann die Nachricht lesen.\\[1mm]
Abgefangener Streifen: \quad\texttt{\large TKLHEIEIECSXAAX}\\[1mm]
Um den Stab passen genau \textbf{3} Buchstaben. Schreibe den Streifen \textbf{spaltenweise} (von oben nach unten) in das Raster und lies \textbf{zeilenweise}.\par
{\lueckentext Klartext: \lsg[6cm]{THE CAKE IS A LIE}\par}
\end{minipage}\hfill
\begin{minipage}[t]{0.38\linewidth}
\vspace{2mm}\centering
\begin{tikzpicture}[x=7mm,y=7mm]
  \draw[gray!70] (0,0) grid (5,3);
  \draw[line width=0.7pt] (0,0) rectangle (5,3);
  \node[font=\small,anchor=east] at (-0.1,2.5) {Zeile 1};
  \node[font=\small,anchor=east] at (-0.1,1.5) {Zeile 2};
  \node[font=\small,anchor=east] at (-0.1,0.5) {Zeile 3};
  \node[font=\ttfamily] at (0.5,2.5) {T}; \node[font=\ttfamily] at (0.5,1.5) {K}; \node[font=\ttfamily] at (0.5,0.5) {L};
  \draw[-{Stealth},blue!60!black,thick] (0.5,3.75) -- (0.5,3.15);
  \ifloesung
  \foreach \z [count=\r] in {{T,H,E,C,A},{K,E,I,S,A},{L,I,E,X,X}} {
    \foreach \b [count=\s] in \z {\ifnum\s>1 \node[font=\ttfamily,text=red!75!black] at (\s-0.5,3.5-\r) {\b};\fi}
  }
  \fi
\end{tikzpicture}\\[1mm]
{\small (X = Füllzeichen)}
\end{minipage}
\end{aufgabe}

\begin{minipage}[t]{0.63\linewidth}
\vspace{0pt}
Helga ist bei einer Recherche zu \textbf{Verschlüsselungen} auf einen Forumseintrag gestoßen. Die Antwort wurde mit dem \textbf{ROT13-Verfahren} verschlüsselt: Jeder Buchstabe des Alphabets wird um 13 Stellen weitergeschoben (nach Z geht es bei A weiter).
\end{minipage}\hfill
\begin{minipage}[t]{0.33\linewidth}
\vspace{0pt}
\begin{tcolorbox}[colback=yellow!15,colframe=orange!60,boxrule=0.5pt,arc=1mm,left=1.5mm,right=1.5mm,top=1mm,bottom=1mm]
\small\itshape Wem wird die Entschlüsselung der Vigenère-Chiffre im 19.\,Jahrhundert zugeschrieben?\\[0.5mm]
\upshape Antwort: \texttt{\bfseries PUNEYRF ONOONTR}
\end{tcolorbox}
\end{minipage}

\begin{aufgabe}[ROT13]{2}
\textbf{a)} Ergänze die Tabelle: Unter jedem Buchstaben steht der Buchstabe, der 13 Stellen weiter liegt. Entschlüssle dann Helgas Antwort.\par\vspace{1.5mm}
{\setlength{\tabcolsep}{0pt}\renewcommand{\arraystretch}{1.25}\small
\begin{tabularx}{\linewidth}{|*{26}{C|}}\hline
A&B&C&D&E&F&G&H&I&J&K&L&M&N&O&P&Q&R&S&T&U&V&W&X&Y&Z\\\hline
\schreibzeile\lsgtext{N}&\lsgtext{O}&\lsgtext{P}&\lsgtext{Q}&\lsgtext{R}&\lsgtext{S}&\lsgtext{T}&\lsgtext{U}&\lsgtext{V}&\lsgtext{W}&\lsgtext{X}&\lsgtext{Y}&\lsgtext{Z}&\lsgtext{A}&\lsgtext{B}&\lsgtext{C}&\lsgtext{D}&\lsgtext{E}&\lsgtext{F}&\lsgtext{G}&\lsgtext{H}&\lsgtext{I}&\lsgtext{J}&\lsgtext{K}&\lsgtext{L}&\lsgtext{M}\\\hline
\end{tabularx}}\par\vspace{1mm}
{\lueckentext Klartext: \lsg[6cm]{CHARLES BABBAGE}\\
\textbf{b)} Öffne das Projekt \texttt{11-5-1 Verschluesselung} (Website) und erklärt euch gegenseitig, wie die Methoden \texttt{rotiere} und \texttt{rot13} arbeiten. Ergänze:\\
Der Computer rechnet mit dem ASCII-Code: \texttt{'A'} hat den Code 65, \texttt{'Z'} den Code \lsg[1.5cm]{90}. In \nr{2} wird die Verschiebung zum Code addiert. Ist das Ergebnis größer als 90, werden in \nr{3} \lsg[1.5cm]{26} abgezogen. In \texttt{rot13} geht eine \lsg[2.5cm]{for-Schleife} den Text Zeichen für Zeichen durch, verschiebt jedes Zeichen um \lsg[1.5cm]{13} und hängt es hinten an den Geheimtext an.\\
\textbf{c)} Verschlüssle den Geheimtext noch einmal mit ROT13. Warum entsteht wieder der Klartext?\par}
\lsglinien{1}{13 + 13 = 26 Stellen: einmal ganz um das Alphabet herum. Ver- und Entschlüsseln sind gleich.}
\end{aufgabe}

\newpage
\textbf{Grundprinzip der Verschlüsselung}\par\vspace{1mm}
\begin{minipage}[c]{0.5\linewidth}
\centering
\begin{tikzpicture}[font=\small,
  txt/.style={draw,ellipse,minimum width=16mm,minimum height=7mm,inner sep=1pt}]
  \node[txt,fill=green!15] (k1) at (0,0) {Klartext};
  \node[txt,fill=orange!20] (g) at (3.05,0) {Geheimtext};
  \node[txt,fill=green!15] (k2) at (6.1,0) {Klartext};
  \draw[-{Stealth},thick] (k1.north east) to[bend left=40] node[above,align=center,font=\scriptsize] {\scalebox{1.6}{\schluessel}\\Verschlüsselung} (g.north west);
  \draw[-{Stealth},thick] (g.north east) to[bend left=40] node[above,align=center,font=\scriptsize] {\scalebox{1.6}{\schluessel}\\Entschlüsselung} (k2.north west);
  \draw[red!70!black] (-0.8,-0.6) -- (-0.8,-0.75) -- (2.1,-0.75) -- (2.1,-0.6); \node[red!70!black,font=\scriptsize] at (0.65,-1) {Sender};
  \draw[violet] (4.0,-0.6) -- (4.0,-0.75) -- (6.9,-0.75) -- (6.9,-0.6); \node[violet,font=\scriptsize] at (5.45,-1) {Empfänger};
\end{tikzpicture}
\end{minipage}\hfill
\begin{minipage}[c]{0.47\linewidth}
{\lueckentext Im Gegensatz zur Codierung (Kapitel 4: Regeln öffentlich) wird bei der \lsg[3.8cm]{Verschlüsselung} eine Nachricht (\textbf{Klartext}) in unlesbare Daten (\textbf{Geheimtext}) umgewandelt, sodass sie vor Unbefugten geschützt ist. Nur wer zusätzlich den \lsg[3cm]{Schlüssel} kennt, kann den Geheimtext wieder entschlüsseln.\par}
\end{minipage}

\medskip
\textbf{Sicherheitsziele einer Verschlüsselung}\par
\begin{tabularx}{\linewidth}{@{}l X@{}}
\textbf{Vertraulichkeit:} & Eine Nachricht kann nur von Befugten gelesen werden.\\
\textbf{Integrität:} & Eine Nachricht kann nicht unbemerkt verändert werden.\\
\textbf{Authentizität:} & Eine Nachricht kann nicht unter falschem Namen verschickt werden.\\
\textbf{Verbindlichkeit:} & Der Absender kann später nicht bestreiten, die Nachricht verfasst zu haben.\\
\end{tabularx}

\begin{aufgabe}[Welches Ziel ist in Gefahr?]{3}
{\lueckentext
\nr{1} Im Schul-WLAN liest jemand deine Chatnachrichten mit: \lsgrest{Vertraulichkeit}
\nr{2} Bei einer Online-Überweisung werden unterwegs aus 10\,€ plötzlich 100\,€: \lsgrest{Integrität}
\nr{3} Eine E-Mail sieht aus, als käme sie von deiner Bank, stammt aber von Betrügern: \lsgrest{Authentizität}
\nr{4} Jemand bestellt online und behauptet später, er habe nie bestellt: \lsgrest{Verbindlichkeit}
\par}
\end{aufgabe}

\textbf{Ein Schlüssel oder ein Schlüsselpaar?}\par
{\lueckentext Man unterscheidet vor allem zwei Arten von Verfahren:\\
\lsg[3.8cm]{Symmetrische} Verschlüsselung: Derselbe Schlüssel wird zum Ver- und Entschlüsseln benutzt.\\
\lsg[3.8cm]{Asymmetrische} Verschlüsselung: Es gibt ein \textbf{Schlüsselpaar}. Mit dem einen Schlüssel kann man nur verschlüsseln, mit dem anderen nur entschlüsseln.\par}

\begin{aufgabe}[The Mystery of Cryptocastle]{4}
Öffne das Spiel (Link auf der Website), gehe zu Abschnitt \textbf{1} und löse die Aufgaben bis zu der Truhe, in der sich der Code \texttt{\bfseries GOVTGYZTTMZKCYE} befindet. Den Code knackst du in Station 5.2.
\end{aufgabe}

\vfill
\web{Station 5.1 Grundlagen: Skytale, Chiffrierscheibe, Code-Animation zu rot13}
\end{document}
;
%            füllt alle Lücken rot aus; siehe \lsg, \lsgrest, \lsglinien,
%            \lsgfeld, \lsgkariert, \lsgkarobox, \lsgtext, \lsgoder)
%  Lösungseinträge belegen nie mehr Platz als die Lücke im Blatt.
% =====================================================================
\usepackage[a4paper,left=18mm,right=18mm,top=26mm,bottom=17mm,headheight=15mm,headsep=4mm,footskip=8mm]{geometry}
\usepackage{iftex}
\ifPDFTeX
  \usepackage[T1]{fontenc}
  \usepackage[utf8]{inputenc}
  \usepackage{helvet}
  \usepackage[scaled=0.86]{beramono}
  \renewcommand{\familydefault}{\sfdefault}
\else
  \usepackage{fontspec}
  \setmainfont{texgyreheros}[Extension=.otf,UprightFont=*-regular,BoldFont=*-bold,ItalicFont=*-italic,BoldItalicFont=*-bolditalic]
  \setmonofont{DejaVuSansMono}[Extension=.ttf,UprightFont=*,BoldFont=*-Bold,ItalicFont=*-Oblique,BoldItalicFont=*-BoldOblique,Scale=0.86]
\fi
\usepackage[ngerman,shorthands=off]{babel}
\usepackage{xcolor,graphicx,tabularx,array,enumitem,fancyhdr,tikz,listings,multicol,calc,etoolbox,ragged2e,colortbl,xspace,pifont}
\usepackage[most]{tcolorbox}
\usetikzlibrary{arrows.meta,positioning,calc,decorations.pathreplacing,shapes.geometric,shapes.symbols,fit}
\usepackage[hidelinks]{hyperref}
\setlength{\parindent}{0pt}
\setlength{\parskip}{3pt}
\setlist{nosep,leftmargin=*}

% ---------- Lösungsschalter ----------
\newif\ifloesung
\ifdefined\mitloesung\loesungtrue\fi
\newcommand{\lsgfarbe}{\color{red!75!black}}

% ---------- Kopf- und Fußzeile (einheitliche Form) ----------
\newcommand{\lehrkraft}{Hau}
\newcommand{\themenbereich}{5. Verschlüsselung}
\newcommand{\abnummer}{}
\pagestyle{fancy}\fancyhf{}
\renewcommand{\headrulewidth}{0.4pt}
\fancyhead[L]{\small Klasse 11\\Datum:}
\fancyhead[C]{\small Themenbereich:\\\themenbereich}
\fancyhead[R]{\small \lehrkraft\ifloesung\\{\lsgfarbe\bfseries Lösung}\fi}
\fancyfoot[C]{\scriptsize edu-mrh.de\,·\,Informatik 11\,·\,Blatt \abnummer\ifloesung\ (Lösung)\fi\,·\,Seite \thepage}
\newcommand{\abtitel}[2]{\renewcommand{\abnummer}{#1}\begin{center}\Large\bfseries #1\quad\underline{#2}\end{center}\vspace{-1mm}}

% ---------- Boxen ----------
\newenvironment{aufgabe}[2][]{\begin{tcolorbox}[enhanced,breakable,colback=white,colframe=black,boxrule=0.7pt,arc=3mm,left=2mm,right=2mm,top=1.5mm,bottom=1.5mm,before skip=2.5mm,after skip=2.5mm]\textbf{Aufgabe #2}\ifstrempty{#1}{}{ \textit{#1}}\par\vspace{0.6mm}}{\end{tcolorbox}}
\newtcolorbox{merke}[1][Merke]{enhanced,breakable,colback=green!5,colframe=green!40!black,boxrule=0.8pt,arc=2mm,left=2mm,right=2mm,top=1.5mm,bottom=1.5mm,before skip=2.5mm,after skip=2.5mm,title={#1},fonttitle=\bfseries,coltitle=white}
\newtcolorbox{hinweis}{enhanced,breakable,colback=yellow!12,colframe=orange!70!black,boxrule=0.6pt,arc=2mm,left=2mm,right=2mm,top=1mm,bottom=1mm,before skip=2mm,after skip=2mm}
\newcommand{\web}[1]{{\small\color{gray}\textit{Website: #1}}}
% Kasten für optionale Aufgaben
\newtcolorbox{optional}[1][Optional]{enhanced,breakable,colback=blue!3,colframe=blue!45!black,boxrule=0.6pt,arc=2mm,left=2mm,right=2mm,top=1.5mm,bottom=1.5mm,before skip=2.5mm,after skip=2.5mm,title={#1},fonttitle=\bfseries\small,coltitle=white}
% Zahl mit Basis: \zs{1011}{2}
\newcommand{\zs}[2]{#1\textsubscript{#2}}

% ---------- Lücken, Linien, Karos ----------
% Einheitliche Maße für alle Blätter der Sequenz:
%  \linien{n}       n Schreiblinien, Abstand 8 mm, erste Linie 8 mm unter dem Text
%  \kariert{n}      Karoraster 5 mm, n Kästchen hoch, nur ganze Kästchen, mittig
%  \karobox{b}{h}   Karoraster b x h, auf ganze Kästchen gerundet
%  \luecke[l]       Lücke im Fließtext (Standard 4,5 cm, Fachbegriffe mind. 3,5 cm)
%  (jeweils mit Lösungsvariante \lsg…, siehe Kopf der Datei)
%  \restzeile       Lücke bis zum Zeilenende
%  \lueckentext     am Anfang eines Absatzes/Kastens mit Lücken: größerer
%                   Zeilenabstand und Flattersatz, damit man in die Lücken schreiben kann
%  \linienfeld{n}   umrahmtes Feld mit n Schreiblinien (z. B. für Java-Code)
%  \schreibzeile    Strut für Tabellenzellen, in die geschrieben wird (8 mm hoch)
\newlength{\zeilenabstand}\setlength{\zeilenabstand}{8mm}
\newlength{\karo}\setlength{\karo}{5mm}
\tikzset{karo/.style={draw=gray!50,line width=0.3pt},schreiblinie/.style={draw=black!65,line width=0.35pt}}
\newlength{\lsgbreite}
% Lücke: \luecke[Breite] / mit Lösung \lsg[Breite]{Text} (zu langer Text wird verkleinert)
\newcommand{\lsg}[2][4.5cm]{\leavevmode\ifloesung\settowidth{\lsgbreite}{\,#2\,}%
  \ifdim\lsgbreite>#1\relax\rlap{\raisebox{0.8pt}{\resizebox{#1}{!}{\lsgfarbe\,#2\,}}}%
  \else\rlap{\raisebox{0.8pt}{\lsgfarbe\,#2}}\fi\fi\rule[-0.6ex]{#1}{0.4pt}\xspace}
\newcommand{\luecke}[1][4.5cm]{\lsg[#1]{}}
% Lücke bis Zeilenende: \restzeile / \lsgrest{Text}
\newcommand{\lsgrest}[1]{\leavevmode\unskip\hspace{1.5mm}\ifloesung\rlap{\raisebox{0.8pt}{\lsgfarbe\,#1}}\fi\leaders\hbox{\rule[-0.6ex]{0.5pt}{0.4pt}}\hfill\null\par}
\newcommand{\restzeile}{\lsgrest{}}
% nur in der Lösung sichtbar (Platz bleibt frei): \lsgtext{…}; Alternativen: \lsgoder{Lösung}{Blatt}
\newcommand{\lsgtext}[1]{\ifloesung\leavevmode{\lsgfarbe#1}\fi}
\newcommand{\lsgoder}[2]{\ifloesung\leavevmode#1\else#2\fi}
\newcommand{\lueckentext}{\linespread{1.5}\selectfont\raggedright}
\newcommand{\schreibzeile}{\rule[-2.8mm]{0pt}{8mm}}
% Schritt-Nummern wie im Code und auf der Website: \nr{1} = ①
\newcommand{\nr}[1]{\ding{\numexpr171+#1\relax}}
% Schreiblinien: \linien{n} / \lsglinien{n}{Text} (Text liegt auf den Linien)
\newcommand{\lsglinien}[2]{\par\vspace{0.5mm}\noindent\begin{tikzpicture}
  \path (0,0) (\linewidth-0.4pt,-#1\zeilenabstand-1.5mm);
  \foreach \i in {1,...,#1}{\draw[schreiblinie] (0.2pt,-\i\zeilenabstand) -- (\linewidth-0.2pt,-\i\zeilenabstand);}
  \ifloesung\node[overlay,anchor=north west,inner sep=0pt] at (1mm,-\zeilenabstand+1pt+5mm) {\parbox[t]{\linewidth-2mm}{\lsgfarbe\raggedright\setlength{\baselineskip}{\zeilenabstand}\rule{0pt}{5mm}#2}};\fi
  \end{tikzpicture}\par}
\newcommand{\linien}[1]{\lsglinien{#1}{}}
% Codefeld mit Linien: \linienfeld{n} / \lsgfeld{n}{Code, Zeilen mit \\ getrennt}
\newcommand{\lsgfeld}[2]{\noindent\begin{tikzpicture}
  \draw[line width=0.4pt] (0.2pt,0) rectangle (\linewidth-0.2pt,-#1\zeilenabstand-3mm);
  \foreach \i in {1,...,#1}{\draw[schreiblinie] (2mm,-\i\zeilenabstand) -- (\linewidth-2mm,-\i\zeilenabstand);}
  \ifloesung\node[overlay,anchor=north west,inner sep=0pt] at (2.5mm,-\zeilenabstand+1pt+5mm) {\parbox[t]{\linewidth-5mm}{\lsgfarbe\ttfamily\small\setlength{\baselineskip}{\zeilenabstand}\rule{0pt}{5mm}#2}};\fi
  \end{tikzpicture}}
\newcommand{\linienfeld}[1]{\lsgfeld{#1}{}}
% Karoraster; Lösungszeichnung als TikZ-Code, Ursprung oben links, y nach unten negativ:
% \kariert{n} / \lsgkariert{n}{TikZ},  \karobox{b}{h} / \lsgkarobox{b}{h}{TikZ}
\newcommand{\karoraster}[3]{\begin{tikzpicture}\draw[karo] (0,0) grid[step=\karo] (#1\karo,#2\karo);
  \ifloesung\begin{scope}[overlay,shift={(0,#2\karo)},red!75!black,text=red!75!black,line width=0.8pt]#3\end{scope}\fi\end{tikzpicture}}
\newcommand{\lsgkariert}[2]{\par\vspace{1.5mm}\noindent\pgfmathtruncatemacro{\karoX}{floor(\linewidth/\karo)}%
  \makebox[\linewidth]{\karoraster{\karoX}{#1}{#2}}\par\vspace{1mm}}
\newcommand{\kariert}[1]{\lsgkariert{#1}{}}
\newcommand{\lsgkarobox}[3]{\pgfmathtruncatemacro{\karoX}{round((#1)/\karo)}\pgfmathtruncatemacro{\karoY}{round((#2)/\karo)}\karoraster{\karoX}{\karoY}{#3}}
\newcommand{\karobox}[2]{\lsgkarobox{#1}{#2}{}}
% Lösungs-Klassenkarte (für Zeichnungen im Karo): \lsgkarte{Name}{Attribute}{Methoden}
\newcommand{\lsgkarte}[3]{{\lsgfarbe\setlength{\arrayrulewidth}{0.6pt}\arrayrulecolor{red!75!black}\begin{tabular}{|l|}\hline\textbf{#1}\\\hline #2\\\hline #3\\\hline\end{tabular}}}
% Lösungs-Objektkarte (in TikZ als Knoteninhalt): \node[objekt] {\lsgobjekt{name : Klasse}{attr = wert\\…}};
\newcommand{\lsgobjekt}[2]{{\lsgfarbe\small\arrayrulecolor{red!75!black}\begin{tabular}{@{}l@{}}\textbf{#1}\\\hline #2\end{tabular}}}
\tikzset{objekt/.style={draw,rounded corners=3mm,inner sep=2mm,anchor=north west}}
\newcommand{\klassenkarte}[3]{\begin{tabular}{|l|}\hline\textbf{#1}\\\hline #2\\\hline #3\\\hline\end{tabular}}
\newcommand{\code}[1]{\texttt{#1}}

% ---------- Java-Listings ----------
\lstset{language=Java,basicstyle=\ttfamily\small,keywordstyle=\bfseries,commentstyle=\color{gray!80!black},stringstyle=\color{black},showstringspaces=false,columns=flexible,keepspaces=true,tabsize=3,frame=single,framerule=0.4pt,rulecolor=\color{black},xleftmargin=2mm,framexleftmargin=1mm,aboveskip=2mm,belowskip=2mm,breaklines=true,
 literate={ä}{{\"a}}1 {ö}{{\"o}}1 {ü}{{\"u}}1 {Ä}{{\"A}}1 {Ö}{{\"O}}1 {Ü}{{\"U}}1 {ß}{{\ss}}1 {–}{{--}}1 {„}{{\glqq}}1 {“}{{\grqq}}1}
\lstnewenvironment{java}[1][]{\lstset{#1}}{}
% Lückencode: größerer Zeilenabstand, Lücken mit (*\luecke[3cm]*) im Code
\lstnewenvironment{javaluecke}[1][]{\lstset{basicstyle=\linespread{1.5}\small\ttfamily,escapeinside={(*}{*)},#1}}{}

\newcolumntype{C}{>{\centering\arraybackslash}X}
\newcommand{\kopfzelle}[1]{\cellcolor{blue!8}\textbf{#1}}
% Schlüssel-Symbol
\newcommand{\schluessel}[1][black]{\tikz[baseline=-0.6ex,#1,line width=0.9pt]{\draw (0,0) circle (0.11); \draw (0.11,0) -- (0.5,0); \draw (0.38,0) -- (0.38,-0.1) (0.46,0) -- (0.46,-0.1);}}
\begin{document}
\abtitel{5.1}{Grundlagen der Verschlüsselung}

\begin{aufgabe}[Die Skytale]{1}
\begin{minipage}[t]{0.58\linewidth}
\vspace{0pt}
Die Spartaner wickelten vor 2500 Jahren einen Lederstreifen um einen Stab (die \textbf{Skytale}) und schrieben ihre Nachricht der Länge nach darauf. Abgewickelt sieht der Streifen wie Buchstabensalat aus. Nur wer einen gleich dicken Stab hat, kann die Nachricht lesen.\\[1mm]
Abgefangener Streifen: \quad\texttt{\large TKLHEIEIECSXAAX}\\[1mm]
Um den Stab passen genau \textbf{3} Buchstaben. Schreibe den Streifen \textbf{spaltenweise} (von oben nach unten) in das Raster und lies \textbf{zeilenweise}.\par
{\lueckentext Klartext: \lsg[6cm]{THE CAKE IS A LIE}\par}
\end{minipage}\hfill
\begin{minipage}[t]{0.38\linewidth}
\vspace{2mm}\centering
\begin{tikzpicture}[x=7mm,y=7mm]
  \draw[gray!70] (0,0) grid (5,3);
  \draw[line width=0.7pt] (0,0) rectangle (5,3);
  \node[font=\small,anchor=east] at (-0.1,2.5) {Zeile 1};
  \node[font=\small,anchor=east] at (-0.1,1.5) {Zeile 2};
  \node[font=\small,anchor=east] at (-0.1,0.5) {Zeile 3};
  \node[font=\ttfamily] at (0.5,2.5) {T}; \node[font=\ttfamily] at (0.5,1.5) {K}; \node[font=\ttfamily] at (0.5,0.5) {L};
  \draw[-{Stealth},blue!60!black,thick] (0.5,3.75) -- (0.5,3.15);
  \ifloesung
  \foreach \z [count=\r] in {{T,H,E,C,A},{K,E,I,S,A},{L,I,E,X,X}} {
    \foreach \b [count=\s] in \z {\ifnum\s>1 \node[font=\ttfamily,text=red!75!black] at (\s-0.5,3.5-\r) {\b};\fi}
  }
  \fi
\end{tikzpicture}\\[1mm]
{\small (X = Füllzeichen)}
\end{minipage}
\end{aufgabe}

\begin{minipage}[t]{0.63\linewidth}
\vspace{0pt}
Helga ist bei einer Recherche zu \textbf{Verschlüsselungen} auf einen Forumseintrag gestoßen. Die Antwort wurde mit dem \textbf{ROT13-Verfahren} verschlüsselt: Jeder Buchstabe des Alphabets wird um 13 Stellen weitergeschoben (nach Z geht es bei A weiter).
\end{minipage}\hfill
\begin{minipage}[t]{0.33\linewidth}
\vspace{0pt}
\begin{tcolorbox}[colback=yellow!15,colframe=orange!60,boxrule=0.5pt,arc=1mm,left=1.5mm,right=1.5mm,top=1mm,bottom=1mm]
\small\itshape Wem wird die Entschlüsselung der Vigenère-Chiffre im 19.\,Jahrhundert zugeschrieben?\\[0.5mm]
\upshape Antwort: \texttt{\bfseries PUNEYRF ONOONTR}
\end{tcolorbox}
\end{minipage}

\begin{aufgabe}[ROT13]{2}
\textbf{a)} Ergänze die Tabelle: Unter jedem Buchstaben steht der Buchstabe, der 13 Stellen weiter liegt. Entschlüssle dann Helgas Antwort.\par\vspace{1.5mm}
{\setlength{\tabcolsep}{0pt}\renewcommand{\arraystretch}{1.25}\small
\begin{tabularx}{\linewidth}{|*{26}{C|}}\hline
A&B&C&D&E&F&G&H&I&J&K&L&M&N&O&P&Q&R&S&T&U&V&W&X&Y&Z\\\hline
\schreibzeile\lsgtext{N}&\lsgtext{O}&\lsgtext{P}&\lsgtext{Q}&\lsgtext{R}&\lsgtext{S}&\lsgtext{T}&\lsgtext{U}&\lsgtext{V}&\lsgtext{W}&\lsgtext{X}&\lsgtext{Y}&\lsgtext{Z}&\lsgtext{A}&\lsgtext{B}&\lsgtext{C}&\lsgtext{D}&\lsgtext{E}&\lsgtext{F}&\lsgtext{G}&\lsgtext{H}&\lsgtext{I}&\lsgtext{J}&\lsgtext{K}&\lsgtext{L}&\lsgtext{M}\\\hline
\end{tabularx}}\par\vspace{1mm}
{\lueckentext Klartext: \lsg[6cm]{CHARLES BABBAGE}\\
\textbf{b)} Öffne das Projekt \texttt{11-5-1 Verschluesselung} (Website) und erklärt euch gegenseitig, wie die Methoden \texttt{rotiere} und \texttt{rot13} arbeiten. Ergänze:\\
Der Computer rechnet mit dem ASCII-Code: \texttt{'A'} hat den Code 65, \texttt{'Z'} den Code \lsg[1.5cm]{90}. In \nr{2} wird die Verschiebung zum Code addiert. Ist das Ergebnis größer als 90, werden in \nr{3} \lsg[1.5cm]{26} abgezogen. In \texttt{rot13} geht eine \lsg[2.5cm]{for-Schleife} den Text Zeichen für Zeichen durch, verschiebt jedes Zeichen um \lsg[1.5cm]{13} und hängt es hinten an den Geheimtext an.\\
\textbf{c)} Verschlüssle den Geheimtext noch einmal mit ROT13. Warum entsteht wieder der Klartext?\par}
\lsglinien{1}{13 + 13 = 26 Stellen: einmal ganz um das Alphabet herum. Ver- und Entschlüsseln sind gleich.}
\end{aufgabe}

\newpage
\textbf{Grundprinzip der Verschlüsselung}\par\vspace{1mm}
\begin{minipage}[c]{0.5\linewidth}
\centering
\begin{tikzpicture}[font=\small,
  txt/.style={draw,ellipse,minimum width=16mm,minimum height=7mm,inner sep=1pt}]
  \node[txt,fill=green!15] (k1) at (0,0) {Klartext};
  \node[txt,fill=orange!20] (g) at (3.05,0) {Geheimtext};
  \node[txt,fill=green!15] (k2) at (6.1,0) {Klartext};
  \draw[-{Stealth},thick] (k1.north east) to[bend left=40] node[above,align=center,font=\scriptsize] {\scalebox{1.6}{\schluessel}\\Verschlüsselung} (g.north west);
  \draw[-{Stealth},thick] (g.north east) to[bend left=40] node[above,align=center,font=\scriptsize] {\scalebox{1.6}{\schluessel}\\Entschlüsselung} (k2.north west);
  \draw[red!70!black] (-0.8,-0.6) -- (-0.8,-0.75) -- (2.1,-0.75) -- (2.1,-0.6); \node[red!70!black,font=\scriptsize] at (0.65,-1) {Sender};
  \draw[violet] (4.0,-0.6) -- (4.0,-0.75) -- (6.9,-0.75) -- (6.9,-0.6); \node[violet,font=\scriptsize] at (5.45,-1) {Empfänger};
\end{tikzpicture}
\end{minipage}\hfill
\begin{minipage}[c]{0.47\linewidth}
{\lueckentext Im Gegensatz zur Codierung (Kapitel 4: Regeln öffentlich) wird bei der \lsg[3.8cm]{Verschlüsselung} eine Nachricht (\textbf{Klartext}) in unlesbare Daten (\textbf{Geheimtext}) umgewandelt, sodass sie vor Unbefugten geschützt ist. Nur wer zusätzlich den \lsg[3cm]{Schlüssel} kennt, kann den Geheimtext wieder entschlüsseln.\par}
\end{minipage}

\medskip
\textbf{Sicherheitsziele einer Verschlüsselung}\par
\begin{tabularx}{\linewidth}{@{}l X@{}}
\textbf{Vertraulichkeit:} & Eine Nachricht kann nur von Befugten gelesen werden.\\
\textbf{Integrität:} & Eine Nachricht kann nicht unbemerkt verändert werden.\\
\textbf{Authentizität:} & Eine Nachricht kann nicht unter falschem Namen verschickt werden.\\
\textbf{Verbindlichkeit:} & Der Absender kann später nicht bestreiten, die Nachricht verfasst zu haben.\\
\end{tabularx}

\begin{aufgabe}[Welches Ziel ist in Gefahr?]{3}
{\lueckentext
\nr{1} Im Schul-WLAN liest jemand deine Chatnachrichten mit: \lsgrest{Vertraulichkeit}
\nr{2} Bei einer Online-Überweisung werden unterwegs aus 10\,€ plötzlich 100\,€: \lsgrest{Integrität}
\nr{3} Eine E-Mail sieht aus, als käme sie von deiner Bank, stammt aber von Betrügern: \lsgrest{Authentizität}
\nr{4} Jemand bestellt online und behauptet später, er habe nie bestellt: \lsgrest{Verbindlichkeit}
\par}
\end{aufgabe}

\textbf{Ein Schlüssel oder ein Schlüsselpaar?}\par
{\lueckentext Man unterscheidet vor allem zwei Arten von Verfahren:\\
\lsg[3.8cm]{Symmetrische} Verschlüsselung: Derselbe Schlüssel wird zum Ver- und Entschlüsseln benutzt.\\
\lsg[3.8cm]{Asymmetrische} Verschlüsselung: Es gibt ein \textbf{Schlüsselpaar}. Mit dem einen Schlüssel kann man nur verschlüsseln, mit dem anderen nur entschlüsseln.\par}

\begin{aufgabe}[The Mystery of Cryptocastle]{4}
Öffne das Spiel (Link auf der Website), gehe zu Abschnitt \textbf{1} und löse die Aufgaben bis zu der Truhe, in der sich der Code \texttt{\bfseries GOVTGYZTTMZKCYE} befindet. Den Code knackst du in Station 5.2.
\end{aufgabe}

\vfill
\web{Station 5.1 Grundlagen: Skytale, Chiffrierscheibe, Code-Animation zu rot13}
\end{document}
